\documentclass[11pt]{article}
\usepackage[margin=.85in]{geometry}
\usepackage{amsmath,amssymb}
\usepackage[colorlinks=true,urlcolor=blue]{hyperref}
\setlength{\parindent}{0pt}
\setlength{\parskip}{.42em}
\setlength{\emergencystretch}{2em}
\newcommand{\ch}{\chi_p}
\newcommand{\leg}[2]{\left(\frac{#1}{#2}\right)}
\newcommand{\F}{\mathbf F_2}
\newcommand{\mollin}{\href{https://www.maths.tcd.ie/EMIS/journals/AMUC/_vol-74/_no_2/_mollin/mollinrea.pdf}{Mollin, \emph{Norm form equations and continued fractions}, (3), (7)--(10)}}
\newcommand{\lls}{\href{https://arxiv.org/pdf/1309.3595}{Lamzouri--Li--Soundararajan, Corollary 1.1}}
\begin{document}
{\Large\bfseries Finding a quadratic nonresidue in a continued-fraction period}\\[3pt]
{\large Individual-prime bounds for \(p\equiv23\pmod{24}\)}

Let \(n_p\) be the least quadratic nonresidue modulo \(p\), \(\ch=(\cdot/p)\),
\(B=\eta(\log p)^2\) with fixed \(\eta>0\), and \(1\le m\le M\ll_\eta(\log\log p)^2\).
We take \(p\) large enough that \(2\sqrt{Mp}<p\) and \(M<B\). Every calculation
of periods below keeps \(p\) fixed. Only the choice of a local multiplier and the
endpoint comparison use Dirichlet GRH.

\section*{1. What occurs in every period}
Write \(D=mp\), let \(a_i/b_i\) be the convergents to \(\sqrt D\) starting at
\(i=0\), and let \(Q_k\) be the positive denominator of the \(k\)th complete
quotient \((P_k+\sqrt D)/Q_k\). For \(1\le k\le\ell_m\), a full period gives
\begin{equation}
 a_{k-1}^2-Db_{k-1}^2=(-1)^k Q_k,\qquad
 0<Q_k<2\sqrt D,\qquad Q_{\ell_m}=1.
\end{equation}
The size bound follows from \(P_k<\sqrt D\) and the complete quotient being
greater than one. Reduction modulo \(p\) gives
\(\ch(Q_k)=(-1)^k\), since \(Q_k<p\) and \(\ch(-1)=-1\). Thus \(\ell_m\)
is even. Each of its \(\ell_m/2\) odd indices has a nonresidue prime factor
to odd exponent. Two is a residue. If any prime divisor of \(m\) is already a
nonresidue, then \(n_p\le M\). For the period witness we can restrict to
multipliers whose prime divisors are all residues, including \(1,2\).
For each such multiplier every negative prime factor above is coprime to \(2mp\).
\mollin.

There is one uniformly small norm in the middle. For \(h=\ell_m/2\),
the symmetry \(P_h=P_{h+1}\) gives \(Q_h\mid2P_h\), and
\(Q_h\mid D-P_h^2\). Hence \(Q_h\mid2D\). As \(Q_h<p\), in fact \(Q_h\mid2m\).
This norm contains only ramified or excluded primes. If all primes dividing
\(m\) are residues, \(\ch(Q_h)=1\) and \(4\mid\ell_m\).
In particular an exceptional prime forces this divisibility for \emph{every}
multiplier in the stated range. \mollin.

For an odd \(q\nmid mp\) dividing \(Q_k\), primitivity of the convergent
gives \(q\nmid b_{k-1}\), hence \((mp/q)=1\). Reciprocity gives
\begin{equation}
 \ch(q)=(-1/q)(p/q)=(-m/q).
\end{equation}
If two multipliers have a common unramified factor \(q\) in their norms, their
predictions agree because \((m_1m_2/q)=1\). Repeated occurrences can still
cancel \(q\) from a product of norms.

Given a prime \(q\le B\) \emph{already satisfying} \(\ch(q)=-1\), choose
\(m=1\) for \(q\equiv3\pmod4\), or \(m=2\) for \(q\equiv5\pmod8\). For
\(q\equiv1\pmod8\), take \(m=2r\) where \(r\) is the least nonresidue mod \(q\).
It is an odd prime and \(r<(\log q)^2\) on GRH. In every case
\((-m/q)=-1\) and \((mp/q)=1\). Taking \(M\ge\max(2,2(\log B)^2)\)
suffices and has the proposed order. The construction starts with \(q\).
Its conclusion is splitting in the auxiliary real quadratic field. \lls.

\newpage
\section*{2. A lower bound that can be evaluated without a candidate}
Choose a nonempty set of multipliers as in section 1 for which every prime
dividing \(m\) is a residue. Take one row \(i=(m,k)\) per period entry
and set \(Q_i=Q_k(mp)\), \(e_i=k\bmod2\). Form the binary matrix
\[
 V_{i,q}=v_q(Q_i)\bmod2 \qquad(q>B\text{ prime dividing some }Q_i).
\]
Let \(W_p(B)\) count distinct primes \(q\le B\) certified by (2) in these
periods, including factors of even multiplicity. Then
\begin{equation}
 W_p(B)\ \ge\
 \operatorname{rank}_{\F}[V\mid e]-\operatorname{rank}_{\F}V .
\end{equation}
This uses only the large factors and the parity of each row index. If the
rank increases, there is a vector \(u\) with \(u^tV=0\) and \(u^te=1\).
The product \(\prod_{u_i=1}Q_i\) has negative character and all prime factors
above \(B\) occur to even exponent. At least one prime \(q\le B\) of odd
exponent in this product has negative character. It divides a selected row
and is coprime to its \(2mp\). Conversely, these two equations are necessary
for a certificate obtained by making the entire large part square and using
the index parities. If the rank increases and there are \(s\) rows, the
number of such subsets is \(2^{s-\operatorname{rank}V-1}\). They can all
certify the same prime. Combining congruences with binary exponent vectors
is classical in continued-fraction factoring.
\href{https://math.dartmouth.edu/~carlp/PDF/paper52.pdf}{Pomerance, \emph{The quadratic sieve factoring algorithm}, section 2}.

Let \(\operatorname{sf}(n)\) be the product of primes occurring to odd
exponent in \(n\), and \(P^+(n)\) its largest prime factor. Then
\begin{equation}
 n_p\ \le\ \min_{\substack{u\in\F^s\\ u^te=1}}
       P^+\!\left(\operatorname{sf}\left(\prod_{u_i=1}Q_i\right)\right).
\end{equation}
The argument of \(P^+\) exceeds one by its character. The right side is the
smallest cutoff giving a rank increase in (3). A long period supplies many
rows. Even if the nullspace of \(V^t\) is large, \(e\) can be orthogonal to
all of it. In fact exceptionality supplies the explicit column relation
\begin{equation}
 n_p>B\quad\Longrightarrow\quad
 e=\sum_{\substack{q>B\\ \ch(q)=-1}} V_{\cdot,q}\quad\text{in }\F^s .
\end{equation}
Here the sum only needs the primes indexing \(V\). Thus using the number
of relations alone drops exactly the parity information needed to exclude
an exceptional prime. A sufficient unproved estimate is a rank increase of
one at \(B=\eta(\log p)^2\) for every sufficiently large \(p\) on the
progression, for each fixed \(\eta>0\). For a single odd row (1) only
puts a factor below \(2\sqrt{mp}\).

At \(p=1{,}478{,}231\), the two periods give
\[
 Q_{15}(p)=235=5\cdot47,\qquad Q_{50}(2p)=1363=29\cdot47.
\]
Their product has character \(-1\) and squarefree part \(145=5\cdot29\).
Since \(p\equiv1\bmod5\), \(5\) is a residue and \(29\) is a nonresidue.
At \(B=(\log p)^2/5=40.3641\ldots\), each individual matrix has rank
difference zero. Stacking them gives rank difference one.

\newpage
\section*{3. Which split primes reach a convergent?}
The local choices \(m=1,2,2r\) from section 1 have squarefree \(D\equiv2,3\bmod4\),
so \(O=\mathbb Z[\sqrt D]\) is the full ring of integers. Take an odd split
prime \(q\nmid D\), choose one ideal \(\mathfrak q\) over \(q\), and define
\[
 b(\mathfrak q)=\min\{N\mathfrak b:\ \mathfrak b\subseteq O
 \text{ a nonzero integral ideal},\
 \mathfrak q\mathfrak b\text{ principal and primitive}\}.
\]
Primitive means the product is not contained in \(tO\) for any integer \(t>1\).
Conjugation gives the same minimum for either prime over \(q\). If \(d_q\)
is the order of \([\mathfrak q]\) in the ordinary class group, \(\mathfrak
q^{d_q}\) is principal and primitive since it lacks the conjugate ideal
factor. In particular \(b(\mathfrak q)\le q^{d_q-1}\). The following two
implications bound the occurrence of \(q\) in this \emph{individual} period.
\begin{equation}
 qb(\mathfrak q)<\sqrt D\quad\Longrightarrow\quad
 q\mid Q_k\text{ for some }k\quad\Longrightarrow\quad
 qb(\mathfrak q)<2\sqrt D .
\end{equation}
For the right implication, the primitive element \(a_{k-1}+b_{k-1}\sqrt D\)
has absolute norm \(Q_k\); its ideal or its conjugate is divisible by
\(\mathfrak q\). For the left one let \(\alpha\) generate the minimizing
product and \(K=|N\alpha|=qb(\mathfrak q)\). Choose its sign so that
\(\alpha>0\) in the positive embedding. Multiply by increasing powers of a
positive Pell unit of norm \(+1\), writing the results \(x+y\sqrt D\).
Primitivity persists, \(x,y>0\) eventually, \(y\to\infty\),
and \(x/y\to\sqrt D\). Thus for \(K<\sqrt D\),
\[
 y^2|x/y-\sqrt D|=\frac{K}{x/y+\sqrt D}<\tfrac12
 \quad\text{eventually}.
\]
Legendre's convergent criterion and periodicity now give \(Q_k=K\)
in the period. In particular \(q^{d_q}<\sqrt D\) suffices for occurrence.
For \(q\asymp B\) this last test asks for
\(d_q<\log D/(2\log q)\sim\log p/(4\log\log p)\). Even a proof of this
order bound for each chosen split prime would still require the existence
of a \(q\le B\) with negative sign. Equation (3) avoids that assumption.

A finite failure of the local inference occurs at \(p=5711\),
\(B=0.26(\log p)^2=19.45\ldots\), \(M=\lfloor(\log\log p)^2\rfloor=4\).
The least nonresidue is \(19\); of \(1\le m\le4\), precisely \(1,4\) pass
the local sign test at \(19\). Their complete denominator lists are
\[
\begin{split}
 &(86,65,43,2,43,65,86,1),\\
 &(43,8,191,109,172,65,235,4,235,65,172,109,191,8,43,1).
\end{split}
\]
Neither contains a multiple of \(19\). There is no period witness at this
cutoff for these four multipliers. This finite example separates splitting
from occurrence. The asymptotic question allows any fixed implied multiplier
constant and a threshold depending on \(\eta\).

\newpage
\section*{4. Finite checks and the endpoint implication}
For \(100{,}000\le p\le2{,}000{,}000\), \(p\equiv23\bmod24\), \(n_p\ge23\),
there are 260 primes. For \(B=(\log p)^2/5\), independent enumeration gives
the following counts. Each successful test counts a prime \(p\) once.
\begin{center}
\begin{tabular}{lr}
Test & Primes\\ \hline
Some \(q\le B\) has \(\ch(q)=-1\), by direct modular exponentiation & 241\\
Some \(q\le B\) certified by (2), \(m=1\) & 198\\
Some \(q\le B\) certified by (2), \(m=1\) or \(2\) & 241\\
Rank increase in (3), \(m=1\) & 197\\
Rank increase in (3), \(m=2\) & 239\\
Rank increase in at least one of the two separate matrices & 240\\
Rank increase in the matrix of both periods together & 241
\end{tabular}
\end{center}
The accompanying \texttt{check\_convergent\_periods.py} generates each period
by the complete-quotient recurrence, checks (1)'s predicted Legendre sign
for every entry, factors the norms, and performs binary elimination. The
rank computation receives neither the direct candidate list nor the signs
of its small prime factors. The unique additional prime obtained by pooling
the two matrices is \(1{,}478{,}231\). The other 19 primes in the cohort
have no nonresidue in the tested range by direct evaluation.

\textbf{The endpoint deficit.} For either parity, put \(Y=(\log p)^2\),
\(E=e^{-\gamma}L(1,\ch)\), and \(A_p=\log Y-E\). The previous endpoint
note's equation (3) gives, on Dirichlet GRH,
\[
 E=\log Y\prod_{\substack{q\le Y\\ \ch(q)=-1}}\frac{q-1}{q+1}+O(1).
\]
The analytic comparison holds for both parities, with \(E\) defined by the
\(L\)-value. All products here run over primes. Set \(n=n_p\); GRH gives
\(n\le Y\) for \(p\ge5\). Inserting every factor from \(n\) to \(Y\) only
decreases the product. Since \((q-1)/(q+1)\ge(1-1/q)^2\), Mertens'
product with the prime number theorem remainder gives
\[
 E+O(1)\ \ge\ \log Y\prod_{n\le q\le Y}(1-1/q)^2
 \ \ge\ \frac{(\log n)^2}{\log Y}-O(1).
\]
For the last inequality use \(\prod_{q<t}(1-1/q)^{-1}
=e^\gamma(\log t+O(1))\), also valid with \(q\le t\).
The lower endpoint error after squaring is \(O(\log n/\log Y)=O(1)\).
Taking square roots and the arithmetic--geometric mean inequality gives
\[
 \log n_p\le\sqrt{\log Y\,(E+O(1))}
 \le \tfrac12(\log Y+E)+O(1),\qquad n_p\ll Y e^{-A_p/2}.
\]
Thus a growing endpoint deficit would suffice, independently of the period
test. \lls. See also Lemmas 2.3 and 2.5 there for the short-product input.

\textbf{Consequence for the growth rate.} Equation (4) is unconditional
on \(23\bmod24\). This pass gives no uniform \(o((\log p)^2)\) bound for
its right side. The \(7\bmod24\) primes have the explicit witness \(3\).
The desired result for all primes, including the \(1\bmod8\) branch,
remains open in this work. The deductions are hand proofs and the interval
check is an integer computation in Python.
\end{document}
